% einstein_photoelectric_ML_slides.tex
% Beamer presentation: "Einstein Was Doing Machine Learning"
% Thomas J. Sargent, 2026
% ~30 minutes, ~22 slides

\documentclass[aspectratio=169,12pt]{beamer}

\usetheme{Madrid}
\usecolortheme{seahorse}
\setbeamertemplate{navigation symbols}{}
\setbeamertemplate{footline}[frame number]

\usepackage{amsmath,amssymb}
\usepackage{booktabs}
\usepackage{array}
\usepackage{multirow}
\usepackage{microtype}
\usepackage{hyperref}
\usepackage{xcolor}

\hypersetup{colorlinks=true, linkcolor=blue, citecolor=blue, urlcolor=blue}

\definecolor{theoryblue}{RGB}{0,80,160}
\definecolor{datared}{RGB}{180,30,30}
\definecolor{mlgreen}{RGB}{30,130,60}

\title[Einstein Was Doing Machine Learning]{%
  \textbf{Einstein Was Doing Machine Learning}\\[4pt]
  {\normalsize The Photoelectric Effect as Statistical Inference}}
\author{Thomas J.\ Sargent}
\institute{New York University \& Hoover Institution}
\date{2026}

%----------------------------------------------------------------------
\begin{document}
%----------------------------------------------------------------------

\begin{frame}[plain]
  \titlepage
\end{frame}

%----------------------------------------------------------------------
\begin{frame}{The Central Claim}

\begin{block}{Thesis}
When Einstein proposed the light-quantum hypothesis in 1905 to explain
the photoelectric effect, he was---in the language we now use---doing
\textbf{machine learning} and \textbf{statistical inference}.
\end{block}

\bigskip
\textbf{Three acts:}
\begin{enumerate}
  \item[\color{theoryblue}1.] \textbf{Theory} (1905): Einstein proposes the model that
    defines the hypothesis class
  \item[\color{datared}2.] \textbf{Data quality} (1887--1916): experimenters spend
    29 years collecting clean training data
  \item[\color{mlgreen}3.] \textbf{Estimation} (1916): Millikan runs the regression
    and recovers Planck's constant to 0.5\%
\end{enumerate}

\bigskip
\small The episode illuminates \emph{both} the power of theory-guided
learning \emph{and} the limits of purely data-driven AI for scientific
discovery.

\end{frame}

%----------------------------------------------------------------------
\begin{frame}{Background: The Photoelectric Effect}

\begin{columns}[T]
\begin{column}{0.55\textwidth}
\textbf{The phenomenon (Hertz, 1887; Hallwachs, 1888):}
\begin{itemize}
  \item Shine light on a metal surface
  \item Electrons are ejected
  \item Apply a \emph{stopping potential} $V_0$ to just halt the
    fastest electrons
\end{itemize}

\bigskip
\textbf{The puzzle (Lenard, 1902):}
\begin{itemize}
  \item[\color{datared}$\bullet$] $V_0$ rises with light \emph{frequency} $\nu$
  \item[\color{datared}$\bullet$] $V_0$ is \textbf{independent} of light
    \emph{intensity} $I$
  \item[\color{datared}$\bullet$] Classical wave theory predicted the opposite
\end{itemize}
\end{column}
\begin{column}{0.42\textwidth}
\begin{block}{Classical Prediction (Wrong)}
Wave amplitude $A \propto \sqrt{I}$ delivers energy\\
$\Rightarrow$ $V_0$ should rise with $I$\\
$\Rightarrow$ $V_0$ should be independent of $\nu$
\end{block}

\bigskip
\begin{alertblock}{What Lenard Found}
Exactly the opposite!\\[4pt]
Intensity $I$: \textbf{irrelevant}\\
Frequency $\nu$: \textbf{everything}
\end{alertblock}
\end{column}
\end{columns}

\end{frame}

%----------------------------------------------------------------------
\begin{frame}[shrink=5]{Lenard's Finding as Feature Engineering}

\begin{columns}[T]
\begin{column}{0.52\textwidth}
\textbf{A modern ML framing of Lenard (1902):}

\bigskip
Lenard ran an ablation experiment.  He measured $V_0$ as a function of
\emph{both} $\nu$ and $I$.

\bigskip
\begin{tabular}{lcc}
\toprule
Feature & Effect on $V_0$ & Importance \\
\midrule
Frequency $\nu$ & Strong $\uparrow$ & \textcolor{mlgreen}{\textbf{High}} \\
Intensity $I$  & None             & \textcolor{datared}{\textbf{Zero}} \\
Wavelength $\lambda$ & (same as $\nu$) & High \\
\bottomrule
\end{tabular}

\bigskip
\textbf{Conclusion:} drop $I$ from the feature set; keep $\nu$.
\end{column}
\begin{column}{0.45\textwidth}
\begin{block}{ML Translation}
\textbf{Input features:} $\nu$, $I$, $\lambda$, \ldots\\[4pt]
\textbf{Label:} stopping potential $V_0$\\[4pt]
\textbf{Lenard's result:} feature-importance
analysis eliminates $I$;
selects $\nu$ as the
\emph{sole relevant input}
\end{block}

\bigskip
``Lenard \ldots established \textbf{two qualitative
facts}: $V_0$ is independent of
intensity; $V_0$ rises with frequency.''

\smallskip
\hfill\small{--- Millikan (1916)}
\end{column}
\end{columns}

\end{frame}

%----------------------------------------------------------------------
\begin{frame}[shrink=5]{Einstein's 1905 Insight: Model Architecture}

\begin{columns}[T]
\begin{column}{0.52\textwidth}
\textbf{The quantum hypothesis:}

\bigskip
Each light quantum carries energy $E = h\nu$.

\smallskip
A single quantum is absorbed by a single electron.

\smallskip
If $h\nu > W$ (work function), the electron escapes.

\bigskip
\textbf{Energy conservation} gives:
\[
  \underbrace{eV_0}_{\text{max KE}}
  = \underbrace{h\nu}_{\text{photon energy}}
  - \underbrace{W}_{\text{work function}}
\]
\end{column}
\begin{column}{0.45\textwidth}
\begin{alertblock}{The Prediction (1905)}
\[
  \boxed{V_0 = \frac{h}{e}\,\nu - \frac{W}{e}}
\]
Linear in frequency.\\[4pt]
Slope $= h/e$ is \textbf{universal} across all metals.\\[4pt]
Intercept $= -W/e$ is metal-specific.
\end{alertblock}

\bigskip
\textbf{Written before a single reliable
$V_0$-vs-$\nu$ measurement existed!}
\end{column}
\end{columns}

\bigskip
\begin{block}{The Inductive Bias}
The quantum hypothesis restricts the hypothesis class from \emph{all
functions} of $(\nu, I, \lambda, \ldots)$ to a \textbf{single affine
function in $\nu$ alone}, with a specific universal slope.
\end{block}

\end{frame}

%----------------------------------------------------------------------
\begin{frame}[shrink=8]{The Regression Equation}

\textbf{Adding measurement noise and a nuisance parameter:}
\begin{equation*}
  V_{0i} = \underbrace{\alpha}_{\text{intercept}} +
            \underbrace{\beta}_{\displaystyle h/e}\,\nu_i +
            \underbrace{\varepsilon_i}_{\text{noise}},
  \qquad
  \alpha = -\frac{W}{e} + \underbrace{v_c}_{\text{contact EMF}}
\end{equation*}

\medskip
\begin{columns}[T]
\begin{column}{0.48\textwidth}
\begin{block}{Parameters of Interest}
\begin{tabular}{ll}
$\beta = h/e$ & slope $\Rightarrow$ Planck's $h$ \\
$\alpha$      & intercept (metal-specific) \\
\end{tabular}
\end{block}

\medskip
\begin{block}{Nuisance Parameter}
$v_c$: contact EMF between emitter
and collector.  Can be as large
as 3\,V---comparable to $V_0$ itself!

\smallskip
Shifts intercept; \textbf{does not bias slope}
if constant within a run.
\end{block}
\end{column}
\begin{column}{0.48\textwidth}
\begin{alertblock}{The Supervised Learning Problem}
\textbf{Training data:} $\{(\nu_i, V_{0i})\}_{i=1}^n$\\[4pt]
\textbf{Task:} estimate $\beta = h/e$\\[4pt]
\textbf{Loss:} squared error\\[4pt]
\textbf{Hypothesis class:} affine functions of $\nu$\\[4pt]
\textbf{Inductive bias:} $\beta > 0$; universal across metals
\end{alertblock}
\end{column}
\end{columns}

\end{frame}

%----------------------------------------------------------------------
\begin{frame}[shrink=10]{29 Years of Collecting Training Data (1887--1916)}

\begin{center}
\small
\begin{tabular}{p{2.8cm}cp{3.8cm}p{4.5cm}}
\toprule
Experimenter & Year & Data quality & Problem \\
\midrule
Hertz / Hallwachs & 1887--88 & Qualitative only & No labeled $(\nu,V_0)$ pairs \\
J.\,J.\ Thomson   & 1899     & Identifies electrons & No $V_0$-vs-$\nu$ measurement \\
Lenard            & 1902     & Qualitative labels & No reliable numerical $V_0$ \\
Ladenburg         & 1907     & 60\,nm range & Linear $\approx$ square-root \\
Hughes            & 1912     & 3 UV lines & Contact EMF uncontrolled; 65\% scatter in $h$ \\
Richardson \& Compton & 1912 & 8 metals  & Same problems; 60\% scatter \\
\midrule
\textbf{Millikan} & \textbf{1916} & \textbf{Clean surfaces, wide $\Delta\nu$, controlled $v_c$} & \textbf{$h$ to 0.5\%} \\
\bottomrule
\end{tabular}
\end{center}

\bigskip
\begin{block}{Data Quality is the Binding Constraint}
The regression is trivial once you have clean data.
The 29-year delay reflects \emph{data engineering}, not statistical
sophistication.
\end{block}

\end{frame}

%----------------------------------------------------------------------
\begin{frame}[shrink=5]{Millikan's Experimental Design (1916)}

\textbf{Key data-quality innovations (``a machine shop in vacuo''):}

\begin{itemize}
  \item[\color{mlgreen}$\checkmark$] \textbf{Fresh surfaces:} electromagnetic knife shaved the alkali-metal cylinder in vacuum---no oxidation, no work-function contamination
  \item[\color{mlgreen}$\checkmark$] \textbf{Spectrally pure lines:} mercury arc with filters to exclude UV companions---eliminates stray-light bias on $V_0$
  \item[\color{mlgreen}$\checkmark$] \textbf{Wide frequency range:} 5461\,\AA\ to 2535\,\AA\
    ($\Delta\nu = 6.3\times10^{14}$\,Hz, nearly a factor of two)
  \item[\color{mlgreen}$\checkmark$] \textbf{Controlled nuisance:} Kelvin double-balance measured $v_c$ continuously; verified stable to $\pm 0.01$\,V over 9 hours
\end{itemize}

\bigskip
\begin{block}{Why the Wide Frequency Range Matters (Cram\'er-Rao)}
\[
  \mathrm{SE}(\hat\beta) = \frac{\sigma}{\sqrt{S_{\nu\nu}}},
  \quad S_{\nu\nu} = \sum_i(\nu_i - \bar\nu)^2
\]
Wider $\Delta\nu$ $\Rightarrow$ larger $S_{\nu\nu}$ $\Rightarrow$ smaller SE.
Hughes's 60\,nm range gives a standard error \textbf{2.2$\times$ larger} than
Millikan's full range.
\end{block}

\end{frame}

%----------------------------------------------------------------------
\begin{frame}[shrink=5]{Millikan's Training Data}

\begin{center}
\renewcommand{\arraystretch}{1.15}
\begin{tabular}{rrcccc}
\toprule
$\lambda$ (\AA) & $\nu$ ($10^{14}$\,Hz)
  & $V_0^{\,\mathrm{Na}}$ (V)
  & $V_0^{\,\mathrm{Li}}$ (V) \\
\midrule
5461.0 &  5.490 & 0.454  & ---   \\
4339.4 &  6.908 & 1.039  & 0.499 \\
4046.8 &  7.409 & 1.246  & 0.706 \\
3650.2 &  8.213 & 1.578  & 1.037 \\
3125.5 &  9.591 & 2.147  & 1.606 \\
2534.7 & 11.827 & 3.030  & 2.529 \\
\bottomrule
\end{tabular}
\end{center}

\medskip
True stopping potential $= V_{0,\text{obs}} + v_c$
\hspace{1em} (Na: $v_c = 2.51$\,V;\; Li: $v_c = 1.52$\,V)

\bigskip
\begin{itemize}
  \item 6 labeled training examples per metal
  \item Two metals (sodium and lithium) available
  \item Points span nearly a \textbf{factor of two} in frequency
\end{itemize}

\end{frame}

%----------------------------------------------------------------------
\begin{frame}[shrink=5]{OLS Regression: Estimating Planck's Constant}

\textbf{Primary regression:} sodium, five points (omitting 2535\,\AA\ scattered-light correction)

\medskip
\begin{columns}[T]
\begin{column}{0.48\textwidth}
\textbf{Standard OLS formula:}
\[
  \hat\beta = \frac{S_{\nu V}}{S_{\nu\nu}}
            = \frac{3.830}{9.277}
            = \mathbf{4.129\times10^{-15}}\,\text{V$\cdot$s}
\]
\[
  \hat\alpha = \bar V_0 - \hat\beta\,\bar\nu = -1.813\,\text{V}
\]

\medskip
\begin{tabular}{lc}
\toprule
Statistic & Value \\
\midrule
$R^2$ & $\approx 0.99999$ \\
Residuals & $\leq 2\,\mathrm{mV}$ \\
SE$(\hat\beta)$ & $4.6\times10^{-18}$\,V$\cdot$s \\
Rel.\ precision & $\mathbf{0.11\%}$ \\
\bottomrule
\end{tabular}
\end{column}
\begin{column}{0.48\textwidth}
\textbf{Implied Planck's constant:}

\medskip
\begin{tabular}{p{3.2cm}r}
\toprule
$e$ used & Implied $h$ \\
\midrule
Millikan's $e$ (1916) & $6.573\times10^{-34}$ \\
Modern NIST $e$ & $6.614\times10^{-34}$ \\
\midrule
\textbf{Modern exact} & $\mathbf{6.626\times10^{-34}}$ \\
\bottomrule
\end{tabular}

\medskip
\begin{alertblock}{Result}
$h$ recovered to \textbf{0.5\%} of its
modern value from \textbf{5 data points}.
\end{alertblock}
\end{column}
\end{columns}

\end{frame}

%----------------------------------------------------------------------
\begin{frame}[shrink=12]{MLE $=$ OLS Under Gaussian Errors}

\textbf{Parametric model:}
\[
  V_{0i} = \alpha + \beta\,\nu_i + \varepsilon_i,
  \quad
  \varepsilon_i \overset{\text{iid}}{\sim} \mathcal{N}(0,\sigma^2)
\]

\textbf{Log-likelihood:}
\[
  \ell(\alpha,\beta,\sigma) =
  -\frac{n}{2}\ln(2\pi) - n\ln\sigma
  - \frac{1}{2\sigma^2}\,\underbrace{\sum_{i}(V_{0i}-\alpha-\beta\nu_i)^2}_{\text{RSS}(\alpha,\beta)}
\]

\bigskip
Maximising $\ell$ over $(\alpha,\beta)$ minimises RSS---exactly OLS.

\[
  \boxed{(\hat\alpha_{\mathrm{MLE}},\hat\beta_{\mathrm{MLE}}) =
         (\hat\alpha_{\mathrm{OLS}},\hat\beta_{\mathrm{OLS}})}
\]

\bigskip
\begin{block}{Implication}
Every OLS result carries over to MLE.
The MSE training loss used in neural-network regression \emph{is} the
Gaussian log-likelihood---the physics and the ML objective coincide.
\end{block}

\end{frame}

%----------------------------------------------------------------------
\begin{frame}[shrink=20]{The Cram\'er-Rao Bound and Experimental Design}

\textbf{Fisher information and Cram\'er-Rao bound for $\beta = h/e$:}
\[
  \mathcal{I}_{\beta\beta} = \frac{S_{\nu\nu}}{\sigma^2},
  \qquad
  \mathrm{Var}(\tilde\beta) \geq \frac{\sigma^2}{S_{\nu\nu}},
  \quad S_{\nu\nu} = \textstyle\sum_i(\nu_i - \bar\nu)^2
\]
OLS achieves this bound (Gauss-Markov + Cram\'er-Rao).

\medskip
\begin{center}
\small
\begin{tabular}{lrrc}
\toprule
Dataset & $n$ & $S_{\nu\nu}$ ($10^{28}$ Hz$^2$) & SE($\hat\beta$) ($10^{-18}$\,V$\cdot$s) \\
\midrule
Restricted (4047--2535\,\AA, like Hughes) & 3 & 1.96  & 10.0 \\
Millikan 5-pt (5461--3126\,\AA)           & 5 & 9.28  & 4.6  \\
Millikan full (5461--2535\,\AA)           & 6 & 28.14 & 2.6  \\
\bottomrule
\end{tabular}
\end{center}

\smallskip
\begin{block}{Design lesson}
Span frequencies widely.  The CR bound justifies Millikan's choice of
a nearly 2$\times$ range in $\nu$ over the narrow 60\,nm UV band used by predecessors.
\end{block}

\end{frame}

%----------------------------------------------------------------------
\begin{frame}[shrink=10]{Universality: Zero-Shot Transfer Across Metals}

\textbf{The quantum hypothesis predicts:} $\beta = h/e$ is the same for
\emph{all} metals.  This is a strong, falsifiable structural restriction.

\bigskip
\begin{columns}[T]
\begin{column}{0.48\textwidth}
\textbf{Separate regressions:}
\begin{align*}
  \hat\beta_{\mathrm{Na}} &= 4.129\times10^{-15}\;\text{V$\cdot$s} \\
  \hat\beta_{\mathrm{Li}} &= 4.137\times10^{-15}\;\text{V$\cdot$s} \\
  \text{difference} &= 0.19\%
\end{align*}
Wald test of $H_0: \beta_{\mathrm{Na}} = \beta_{\mathrm{Li}}$:
\[
  z \approx 0.87 \quad (\ll 1.96)
\]
Cannot reject universality.
\end{column}
\begin{column}{0.48\textwidth}
\begin{alertblock}{ML Language}
\textbf{Pre-training:} fit slope on sodium\\[4pt]
\textbf{Zero-shot transfer:} apply slope
to lithium (different intercept only)\\[4pt]
\textbf{Fine-tuning needed:} only intercept
$\alpha_{\mathrm{Li}}$ (work function shifts)\\[4pt]
\textbf{Transfer error:} 0.2\%
\end{alertblock}
\end{column}
\end{columns}

\bigskip
The pooled estimate $\hat\beta_{\mathrm{pool}} \approx 4.133\times10^{-15}$\,V$\cdot$s
has SE reduced by $\approx 1/\sqrt{2}$ relative to either single-metal estimate.

\end{frame}

%----------------------------------------------------------------------
\begin{frame}[shrink=5]{Historical Estimates of Planck's Constant}

\begin{center}
\begin{tabular}{p{4.4cm}rcr}
\toprule
Authors & Year & $h$ ($10^{-27}$\,erg$\cdot$s) & Error \\
\midrule
Hughes              & 1912 & 3.55--5.85        & $-12$ to $-46\%$ \\
Richardson \& Compton & 1912 & 3.5--5.8        & $-12$ to $-47\%$ \\
Planck (black body) & 1913 & 6.415             & $-3.2\%$         \\
\midrule
Millikan (Na, 5-pt) & 1916 & 6.573             & $-0.80\%$        \\
Millikan (Na, 9-pt) & 1916 & 6.578             & $-0.73\%$        \\
Millikan (Li)       & 1916 & 6.584             & $-0.64\%$        \\
\textbf{Millikan mean} & \textbf{1916} & \textbf{6.579} & $\mathbf{-0.72\%}$ \\
\midrule
Modern NIST 2022    & 2022 & 6.6261            & 0                \\
\bottomrule
\end{tabular}
\end{center}

\bigskip
\textbf{Pattern:} pre-Millikan estimates biased downward 10--47\%.

\smallskip
Three causes: stray UV (inflates $V_0$ at long wavelengths $\Rightarrow$ flattens slope),
narrow $\Delta\nu$ (large SE), uncontrolled $v_c$ (slope bias).

\smallskip
Millikan solved all three.

\end{frame}

%----------------------------------------------------------------------
\begin{frame}[shrink=12]{Millikan's Epistemic Reluctance}

\begin{columns}[T]
\begin{column}{0.52\textwidth}
Millikan confirmed Einstein's equation to 0.5\%.  \textbf{Then wrote:}

\bigskip
\begin{quote}
``Despite then the apparently complete success of the Einstein equation,
the physical theory of which it was designed to be the symbolic expression
is found \textbf{so untenable} that Einstein himself, I believe, no longer
holds to it.''
\end{quote}
\hfill\small{--- Millikan (1916), p.\ 388}

\bigskip
He had hoped the experiment would \textbf{refute} the equation!
\end{column}
\begin{column}{0.45\textwidth}
\begin{block}{Lesson for Statistics}
Confirming a regression equation is \textbf{not} the same as confirming
the structural model that generated it.

\bigskip
Millikan confirmed the slope $h/e$ to 0.5\% without accepting that light
consists of quanta.
\end{block}

\medskip
Compton scattering (1923) finally convinced him.

\smallskip
Nobel acceptance (1923):\\
``\ldots contrary to my own expectation,
\textbf{first direct experimental proof}
\ldots of the complete validity of the
Einstein equation.''
\end{column}
\end{columns}

\end{frame}

%----------------------------------------------------------------------
\begin{frame}[shrink=30]{The ML-to-Physics Dictionary}

\begin{center}
\footnotesize
\renewcommand{\arraystretch}{1.10}
\begin{tabular}{p{4.8cm}p{4.8cm}p{3.5cm}}
\toprule
\textbf{Scientific step} & \textbf{ML name} & \textbf{Who} \\
\midrule
Propose $E = h\nu$; derive $V_0 = (h/e)\nu - W/e$ & Inductive bias + Hypothesis class & Einstein (1905) \\
$V_0$ depends on $\nu$ not $I$ & Feature engineering & Lenard (1902) \\
Measure $(\nu_i, V_{0i})$ pairs & Labeled training data & Millikan (1916) \\
Control contact EMF $v_c$ & Nuisance-variable identification & Millikan (1916) \\
Min.\ $\sum(V_{0i} - \hat V_{0i})^2$ & MSE training (OLS/MLE) & Millikan; this paper \\
$|\hat\beta - h/e| < 1\%$ & Test accuracy / Confirmation & Millikan (1916) \\
Na and Li give same $\hat\beta$ & Zero-shot transfer learning & Millikan (1916) \\
\bottomrule
\end{tabular}
\end{center}

\end{frame}

%----------------------------------------------------------------------
\begin{frame}[shrink=12]{The Hypothesis Class: Where Einstein Differs from a Naive ML System}

\begin{columns}[T]
\begin{column}{0.52\textwidth}
\textbf{Einstein's hypothesis class:}
\[
  \mathcal{H} = \bigl\{f_{\alpha,\beta}: \nu \mapsto \alpha + \beta\nu
  \bigm| \alpha\in\mathbb{R},\;\beta>0\bigr\}
\]
with the \emph{cross-metal restriction} that $\beta$ is universal.

\medskip
\begin{itemize}
  \item VC dimension $= 2$
  \item One-parameter family of parallel lines
  \item Derived \emph{deductively} from $E = h\nu$
  \item Five training points are \emph{more than sufficient}
\end{itemize}
\end{column}
\begin{column}{0.45\textwidth}
\begin{alertblock}{What a Naive ML System Misses}
A linear regression package gives the
right numerical answer---but:

\begin{enumerate}
  \item Cannot identify $\hat\beta = 4.129\times10^{-15}$
    as $h/e$
  \item Cannot impose $\beta_{\mathrm{Na}} = \beta_{\mathrm{Li}}$
    \emph{from first principles}
  \item Cannot decompose intercept into $W/e$ vs.\ $v_c$
  \item Cannot reduce data requirement from $O(10^3)$ to~5
\end{enumerate}
\end{alertblock}
\end{column}
\end{columns}

\bigskip
\begin{block}{The Reduction Einstein Achieved}
$E = h\nu$ reduced the hypothesis space from the infinite-dimensional
function class to a \textbf{two-dimensional manifold of parallel lines}---
compressing the data requirement by three orders of magnitude.
\end{block}

\end{frame}

%----------------------------------------------------------------------
\begin{frame}[shrink=12]{Einstein vs.\ a Deep Neural Network: Three Epochs}

\begin{center}
\small
\renewcommand{\arraystretch}{1.2}
\begin{tabular}{p{2.8cm}rp{2.8cm}p{2.8cm}p{3.5cm}}
\toprule
Epoch & $n$ avail. & Einstein model & DNN (VC $\approx$ 5700) & DNN outcome \\
\midrule
\textbf{Lenard 1902}   & $\approx 0$  & Sufficient$^\dagger$ & Insufficient
  & Cannot train (no labels) \\[2pt]
\textbf{Hughes 1912}   & $\sim 30$    & Sufficient
  & Insufficient ($\times 190$) & Memorises; no generalisation \\[2pt]
\textbf{Millikan 1916} & 5            & Sufficient
  & Catastrophically insuf. & Interpolates; no inference \\[2pt]
\textbf{Modern era}    & $\sim 5000$  & Sufficient
  & Sufficient & Recovers line; cannot identify $h/e$ \\
\bottomrule
\multicolumn{5}{p{14.5cm}}{\small $^\dagger$Functional form derived from theory; only two free parameters remain.}
\end{tabular}
\end{center}

\bigskip
\small A three-layer DNN (16 neurons/layer) has $p \approx 593$ parameters and
VC dim $\approx 5700$.  Reliable 1\% generalisation requires $n \gtrsim 5700$---short by
three orders of magnitude at Millikan's epoch.

\end{frame}

%----------------------------------------------------------------------
\begin{frame}[shrink=8]{The DNN on Millikan's Five Points}

\textbf{What happens if you train any DNN on $n = 5$ data points?}

\bigskip
\begin{columns}[T]
\begin{column}{0.48\textwidth}
\textbf{Interpolation regime} (any DNN with $p > 5$):
\begin{itemize}
  \item Training loss $= 0$ (memorises all points)
  \item Function is unconstrained \emph{between} points
  \item Cannot estimate slope uncertainty: zero residual df
\end{itemize}

\bigskip
\textbf{OLS with $n=5$, $p=2$:}
\begin{itemize}
  \item 3 residual degrees of freedom
  \item $\mathrm{SE}(\hat\beta) = 4.6\times10^{-18}$\,V$\cdot$s
  \item 0.11\% relative precision, at the Cram\'er-Rao bound
\end{itemize}
\end{column}
\begin{column}{0.48\textwidth}
\begin{alertblock}{The Core Lesson}
\textbf{Einstein's theory} (not more data,
not a better algorithm) is what made
five training points sufficient.

\bigskip
Specifying the hypothesis class was a
harder intellectual act than fitting
the line.
\end{alertblock}

\bigskip
\begin{block}{Where DNNs Win}
High-dimensional problems where no
parsimonious parametric family is known
(images, proteins, language).

\smallskip
The photoelectric problem sits at the
\emph{opposite} extreme.
\end{block}
\end{column}
\end{columns}

\end{frame}

%----------------------------------------------------------------------
\begin{frame}[shrink=5]{A Broader Lesson: The Hypothetico-Deductive Loop}

\begin{center}
\small
\renewcommand{\arraystretch}{1.25}
\begin{tabular}{p{4.0cm}p{4.0cm}p{4.0cm}}
\toprule
\textbf{Step} & \textbf{Science} & \textbf{ML} \\
\midrule
Propose $E = h\nu$              & Theoretical hypothesis & Model architecture \\
Derive $V_0 = (h/e)\nu - W/e$   & Deduction              & Hypothesis class \\
Choose $\nu$, not $I$           & Experiment design      & Feature engineering \\
Measure $(\nu_i, V_{0i})$       & Observation            & Training data \\
Minimise RSS                    & Least squares          & MSE training \\
$|\hat\beta - h/e| < 1\%$       & Confirmation           & Test accuracy \\
Na and Li: same $\hat\beta$     & Universality test      & Zero-shot transfer \\
\bottomrule
\end{tabular}
\end{center}

\bigskip
\begin{block}{Modern AI for Science --- the Missing Piece}
Modern systems (AI Feynman, symbolic regression) can \emph{induce} the
linear law from data by exhaustive search, but they do so over a
combinatorially large function space.  Einstein's deductive step
encoded the result of that search in one paragraph.

\smallskip
\textbf{The human contribution---theory construction---remains
indispensable.}
\end{block}

\end{frame}

%----------------------------------------------------------------------
\begin{frame}[shrink=5]{Bonus: Born's Rule and an Unintended Consequence}

\begin{columns}[T]
\begin{column}{0.52\textwidth}
Einstein's quantum hypothesis required:
\[
  |E|^2 \;\propto\; \text{photon \emph{density}}
\]
The classical wave intensity $|E|^2$ becomes a
\textbf{probability density}, not just an energy density.

\medskip
Einstein called this the \emph{Gespensterfeld}
(ghost field)---a wave that guides probabilities.

\bigskip
\textbf{Max Born (1926)} formalised the same logic
for the Schr\"odinger wave function:
\[
  \boxed{|\psi|^2 = \text{probability density}}
\]
\textbf{Born's rule} --- now the cornerstone of quantum mechanics.
\end{column}
\begin{column}{0.45\textwidth}
\begin{alertblock}{Two Nobel Prizes}
\textbf{1921:} Einstein --- for the equation
that made the measurement possible
\\[4pt]
\textbf{1954:} Born --- for the probabilistic
interpretation the equation made
\emph{inevitable}
\end{alertblock}

\bigskip
Einstein spent the rest of his life
rejecting Born's rule as a ``complete
description of reality.''

\bigskip
He won the Nobel for a paper whose
logic seeded the interpretation he
most opposed.
\end{column}
\end{columns}

\end{frame}

%----------------------------------------------------------------------
\begin{frame}[shrink=5]{Summary and Conclusions}

\begin{block}{The Thesis, Revisited}
Einstein was doing machine learning in 1905---but the most important step
was the one no ML system has yet automated: \textbf{specifying the
hypothesis class}.
\end{block}

\bigskip
\textbf{Three principles illustrated:}
\begin{enumerate}
  \item \textbf{Theoretical priors guide design.} Once $E=h\nu$ fixes the
    functional form, only 5 training points are needed for 0.1\% precision.
  \item \textbf{Data quality is the binding constraint.} The regression is
    trivial; cleaning the data took 29 years.
  \item \textbf{Confirming an equation $\neq$ endorsing the model.}
    Millikan confirmed the slope to 0.5\% while publically rejecting the
    light-quantum interpretation.
\end{enumerate}

\bigskip
\begin{alertblock}{The Limit of Pure AI}
A data-driven system given only $(\nu_i, V_{0i})$ pairs can recover
$\hat\beta = 4.129\times10^{-15}$\,V$\cdot$s but cannot know it is $h/e$.
\textbf{Physical meaning requires physical theory.}
\end{alertblock}

\end{frame}

%----------------------------------------------------------------------
\begin{frame}[allowframebreaks]{References}

\small
\begin{itemize}
  \item Einstein, A.\ (1905). ``\"Uber einen die Erzeugung und Verwandlung des
    Lichtes betreffenden heuristischen Gesichtspunkt.''
    \textit{Ann.\ Phys.}, 17, 132--148.
  \item Lenard, P.\ (1902). ``\"Uber die lichtelektrische Wirkung.''
    \textit{Ann.\ Phys.}, 8, 149--198.
  \item Millikan, R.\,A.\ (1916). ``A Direct Photoelectric Determination of
    Planck's $h$.''
    \textit{Phys.\ Rev.}, 7(3), 355--388.
  \item Born, M.\ (1926). ``Zur Quantenmechanik der Sto{\ss}vorg\"ange.''
    \textit{Z.\ Phys.}, 37, 863--867.
  \item Hughes, A.\,L.\ (1913). ``On the Emission Velocities of
    Photo-Electrons.'' \textit{Phil.\ Trans.\ R.\ Soc.\ A}, 212, 205--226.
  \item Richardson, O.\,W.\ \& Compton, K.\,T.\ (1912). ``The
    Photoelectric Effect.'' \textit{Phil.\ Mag.}, 24, 575--594.
  \item Udrescu, S.\,M.\ \& Tegmark, M.\ (2020). ``AI Feynman: A
    Physics-Inspired Method for Symbolic Regression.''
    \textit{Sci.\ Adv.}, 6, eaay2631.
  \item Vapnik, V.\ (1998). \textit{Statistical Learning Theory}.
    Wiley, New York.
  \item Hornik, K.\ et al.\ (1989). ``Multilayer Feedforward Networks Are
    Universal Approximators.'' \textit{Neural Networks}, 2(5), 359--366.
  \item Belkin, M.\ et al.\ (2019). ``Reconciling Modern Machine-Learning
    Practice and the Classical Bias-Variance Trade-Off.''
    \textit{PNAS}, 116(32), 15849--15854.
\end{itemize}

\end{frame}

%----------------------------------------------------------------------
\end{document}
